% -----------------------------------------------------------
\section{Firstborn}
% -----------------------------------------------------------
	\label{FIRSTBORN}
	
% % ----------------------
% \subsection{See also}
% % ----------------------

% % ----------------------
% \subsection{1828 American Dictionary of the English Language} % *1828
% % ----------------------
	% \label{FIRSTBORN-1828}

	% \begin{compactitem}
		% \item
	% \end{compactitem}

% % ----------------------
% \subsection{Oxford English Dictionary} % *OED
% % ----------------------
	% \label{FIRSTBORN-OED}

	% \begin{compactitem}
		% \item[]
	% \end{compactitem}
	
% % ----------------------
% \subsection{Restoration Lexicon} % *RL *RESTORATION LEXICON
% % ----------------------
	% \label{FIRSTBORN-OED}

	% \begin{compactitem}
		% \item[]
	% \end{compactitem}

% ----------------------
\subsection{Scriptural Usage} % *SCRIPTURES
% ----------------------
	\label{FIRSTBORN-SCRIPTURES}

	% % -----
	% \subsubsection{Old Testament} % *OT
	% % -----
		% \label{FIRSTBORN-OT}
	
		% % --
		% \paragraph{KJV}
		% % --
			% \label{FIRSTBORN-OT-KJV}
	
			% \begin{tabular}{ll}
				% Total occurrences in the OT & 
			% \end{tabular}
			
			% \begin{compactdesc}
				% \item[]
			% \end{compactdesc}
			
		% % --
		% \paragraph{Hebrew Bible}
		% % --
			% \label{FIRSTBORN-HEBREW}
		
			% %
			% \subparagraph{\texorpdfstring{\he{sample word}}{}} % paste actual hebrew text here
			% %
				% \label{PAGA} % whatever is in step bible without periods
			
				% \begin{compactdesc}
					% \item[HALOT , PDF ] \textbf{} % Use the 1994 PDF
						% \begin{compactitem}
							% \item[1]
						% \end{compactitem}
					% \item[BDB , PDF ] \textbf{}
						% \begin{compactitem}
							% \item[1]
						% \end{compactitem}
					% \item[DCH PDF ] \textbf{} % don't include the actual page number, they repeat from volume to volume
						% \begin{compactitem}
							% \item[1]
						% \end{compactitem}
				% \end{compactdesc}
				
				% \begin{compactdesc}
					% \item[Some OT uses] \textbf{}
						% \begin{compactdesc}
							% \item[]
						% \end{compactdesc}
				% \end{compactdesc}
			
		% % --
		% \paragraph{Septuagint}
		% % --
			% \label{FIRSTBORN-LXX}
		
			% %
			% \subparagraph{\texorpdfstring{\gr{sample word}}{}}
			% %
		
	% -----
	\subsubsection{New Testament} % *NT
	% -----
		\label{FIRSTBORN-NT}
	
		% % --
		% \paragraph{KJV}
		% % --
			% \label{FIRSTBORN-NT-KJV}		
	
			% \begin{tabular}{ll}
				% Total occurrences in the NT & 
			% \end{tabular}
			
			% \begin{compactdesc}
				% \item[]
			% \end{compactdesc}
			
		% --
		\paragraph{Greek New Testament} % *GREEK
		% --
			\label{FIRSTBORN-GREEK}
		
			%
			\subparagraph{\texorpdfstring{\gr{πρωτότοκος}}{}}
			%
				\label{PROTOTOKOS}
			
				\begin{compactdesc}
					\item[BDAG 793a] \textbf{}
						\begin{compactitem}
							\item[1] literally \textbf{pertaining to birth order, \textit{firstborn}}...the special status enjoyed by a firstborn son as heir apparent in Israel is an implicit component of \gr{πρωτότοκος} in reference to such a son and plays a dominant role in [definition \#2]
							\item[2] \textbf{pertaining to having special status associated with a firstborn, \textit{firstborn}}, figurative extension of definition \#1
							\item[2.a] of Christ, as the firstborn of a new humanity which is to be glorified, as its exalted Lord is glorified \gr{πρωτότοκος ἐν πολλοῖς ἀδελφοῖς} (\hyperref[ROMANS-8-29]{Romans 8:29}). Also simply \gr{πρωτότοκος} (\hyperref[HEBREWS-1-6]{Hebrews 1:6}); compare \hyperref[REVELATIONNT-2-1]{Revelation 2:1}. This expression, which is admirably suited to describe Jesus as the one coming forth from God to found the new community of believers, is also used in some instances where the force of the element -\gr{τοκος} appears at first glance to be uncertain...
							\item[2.b] of ordinary humans
							\item[2.b.$\alpha$] of God's people...of the \textit{assembly of the firstborn} (see \hyperref[EXODUS-4-22]{Exodus 4:22}) in heaven \gr{ἐκκλησία πρωτοτόκων} (\hyperref[HEBREWS-12-23]{Hebrews 12:23})
							\item[2.b.$\beta$] of a notorious dissident \gr{πρωτότοκος τοῦ Σατανᾶ}
								\begin{compactitem}
									\item This last one is some special use in certain Christian writers, it has no bible references associated with it.
								\end{compactitem}
						\end{compactitem}
					\item[CGL 1235b, PDF 1275] \textbf{}
						\begin{compactitem}
							\item (of a woman, cow, goat) \textbf{that has given birth for the first time}; (of a thinker, figuratively referring to having produced a doctrine) (Plato)
							\item ---\gr{πρωτότοκος} (of a child) \textbf{first-born} (NT)
						\end{compactitem}
					\item[LSJ 1545b, PDF 1616] \textbf{}
						\begin{compactitem}
							\item[I] \textit{bearing} or \textit{having borne her first-born}
							\item[II.1] \textit{first-born}
							\item[II.2] metaphorical, \gr{πρωτότοκος πάσης κτίσεως} (\hyperref[COLOSSIANS-1-15]{Colossians 1:15})
						\end{compactitem}
				\end{compactdesc}
				
				% \begin{tabular}{ll}
					% Total occurrences in the NT & 
				% \end{tabular}
				
				% \begin{compactdesc}
					% \item[Some NT uses] \textbf{}
						% \begin{compactdesc}
							% \item[]
						% \end{compactdesc}
				% \end{compactdesc}
				
			% %
			% \subparagraph{TDNT: \texorpdfstring{\gr{sample word}}{}  (3:536--556)}
			% %
				% \label{KATALLAGE-TDNT}
		
	% % -----
	% \subsubsection{Book of Mormon} % *BOM
	% % -----
		% \label{FIRSTBORN-BOM}
	
		% \begin{tabular}{ll}
			% Total occurrences in the BOM & 
		% \end{tabular}
		
		% \begin{compactdesc}
			% \item[]
		% \end{compactdesc}
		
	% % -----
	% \subsubsection{Doctrine and Covenants} % *DC
	% % -----
		% \label{FIRSTBORN-DC}
	
		% \begin{tabular}{ll}
			% Total occurrences in the DC & 
		% \end{tabular}
		
		% \begin{compactdesc}
			% \item[]
		% \end{compactdesc}
		
	% % -----
	% \subsubsection{Pearl of Great Price} % *PGP
	% % -----
		% \label{FIRSTBORN-PGP}
	
		% \begin{tabular}{ll}
			% Total occurrences in the PGP & 
		% \end{tabular}
		
		% \begin{compactdesc}
			% \item[]
		% \end{compactdesc}
		
% ----------------------
\subsection{Key Phrases} % *KEYPHRASES *PHRASES
% ----------------------
	\label{FIRSTBORN-PHRASES}

	\begin{compactdesc}
		\item[church of the firstborn] See this phrase at \hyperref[CHURCH-PHRASES]{here}.
			% \begin{compactdesc}
				% \item[Bible anchor verses] \phantom{.}
					% \begin{compactdesc}
						% \item[Romans 5:5] \gr{}
					% \end{compactdesc}
				% \item[Commentary] ---
			% \end{compactdesc}
	\end{compactdesc}
	
% % ----------------------
% \subsection{Bible Dictionaries} % *DICTIONARIES
% % ----------------------
	% \label{FIRSTBORN-BD}

% % ----------------------
% \subsection{Restoration Usage} % *RESTORATION
% % ----------------------
	% \label{FIRSTBORN-RESTORATION}

	% % -----
	% \subsubsection{Joseph Smith} % *JS
	% % -----
		% \label{FIRSTBORN-JS}
	
		% \begin{compactdesc}
			% \item[{\hyperref[KEY]{Date/Source}}]
		% \end{compactdesc}
		
	% % -----
	% \subsubsection{Temple Ordinances}
	% % -----
		% \label{FIRSTBORN-TEMPLE}
	
		% \begin{compactdesc}
			% \item[{\hyperref[KEY]{Date/Source}}]
		% \end{compactdesc}
	
	% % -----
	% \subsubsection{Brigham Young} % *BY
	% % -----
		% \label{FIRSTBORN-BY}
	
		% \begin{compactdesc}
			% \item[{\hyperref[KEY]{Date/Source}}]
		% \end{compactdesc}
	
% % ----------------------
% \subsection{Commentary and Studies} % *COMMENTARY *STUDIES
% % ----------------------
	% \label{FIRSTBORN-COMMENTARY}

	% % -----
	% \subsubsection{Key Points}
	% % -----
		% \label{FIRSTBORN-KEYS}
	
		% \begin{compactitem}
			% \item
		% \end{compactitem}
		
	% % -----
	% \subsubsection{Sample Study}
	% % -----
		% \label{FIRSTBORN-study}